%%%%%%%%%%% LANGUAGE %%%%%%%%%%%

% International language support
% Fetches language from documentclass options. Most other packages do this as well
\usepackage{babel}

% Font encoding
% Without this package, words containing åäö are not hyphenated automatically
\usepackage[T1]{fontenc}

% For interpreting non-ASCII characters
\usepackage[utf8]{inputenc}


%%%%%%%%%%% FORMAL STUFF %%%%%%%%%%%

% Smaller margins
% \usepackage[margin=2.5cm]{geometry}

% % Clickable urls
% \usepackage[hyphens]{url} % Does not want to be loaded after physics package

% Dates & time
\usepackage[yyyymmdd]{datetime} % Useful when referencing websites
\renewcommand{\dateseparator}{-} % ISO 8601 date format

% What to display in table of contents
\setcounter{tocdepth}{1}
\setcounter{secnumdepth}{2}

% Lists
\usepackage{enumerate} % Determines the style in which the counter is printed
\usepackage{enumitem} % Provides user control over the layout of the three basic list environments

% Citing & bibliography
\usepackage{csquotes} % For \enquote command for proper quotation marks, biblatex uses this if it is loaded
\usepackage[style=alphabetic]{biblatex} 
\bibliography{bibliography} % If not used, the preceding two lines should be commented out. Otherwise LaTeX will complain of an empty bibliography.

% Clickable links and refs
\usepackage[hidelinks]{hyperref} 


%%%%%%%%%%% GRAPHICS %%%%%%%%%%%

\usepackage{color}
\usepackage{xcolor}

% Figures
\usepackage{graphics}
\usepackage{graphicx} % More options for \includegraphics
% \usepackage{epsfig} % Solves some problems in \includegraphics{<.eps-file>}
\usepackage{float} % More and improved options for figure placement
\usepackage{caption} % More options for \caption
\usepackage{subcaption} % Subfigures
\usepackage{adjustbox} % \adjustbox{valign=c}{...} useful for vertically aligning graphics in equations

% % PGF/TikZ
% % They are two "low level" languages for producing graphics
% % I often create graphics in a separate file using the 'standalone' documentclass
% \usepackage{pgf}
% \usepackage{pgfplots}
% \pgfplotsset{compat=1.18}
% \usepackage{tikz}
% \usetikzlibrary{calc} % Do calculations in coordinate expressions
% \usetikzlibrary{graphs} % Useful for Penrose diagrams, commutative diagrams, etc
% \usetikzlibrary{quotes} % Edge labels
% \usetikzlibrary{arrows.meta} % Arrow head scaling, etc

% For the Aldus leaf (\ding{166})
\usepackage{pifont}


%%%%%%%%%%% PHYSICS %%%%%%%%%%%

% SI units
\usepackage{siunitx}[=v2] % v2 option to avoid warning about \qty
\DeclareSIUnit\clight{\text{$c$}} % redefine from c_0 to c
\DeclareSIUnit\byte{B}

% Physics macros
\usepackage[notrig]{physics} % Defines lots of nice commands like \derivative, \norm, \evaluated, etc. If loaded without notrig option, it redefines \sin, \cos, etc. This is a strange default...
\usepackage{braket} % Defines \bra, \ket, \braket, and \set
\usepackage{slashed} % For Feynman slash notation
% \usepackage{simpler-wick} % Wick contractions (may require sty-file)
% \usepackage[compat=1.1.0]{tikz-feynman} % Feynman diagrams (has to be compiled with LuaTeX)
% \usepackage{tensor} % Covariant index notation (I don't use this anymore. I just do T^i{}_j{}_k)


%%%%%%%%%%% CODING %%%%%%%%%%%

% For nice code insertions
\usepackage{listings}
\definecolor{codegreen}{rgb}{0,0.6,0}
\definecolor{codegray}{rgb}{0.5,0.5,0.5}
\definecolor{codepurple}{rgb}{0.58,0,0.82}
\definecolor{backcolour}{rgb}{0.95,0.95,0.92}
\lstdefinestyle{mystyle}{
    backgroundcolor=\color{backcolour},   
    commentstyle=\color{codegreen},
    keywordstyle=\color{magenta},
    numberstyle=\tiny\color{codegray},
    stringstyle=\color{codepurple},
    basicstyle=\ttfamily\footnotesize,
    breakatwhitespace=false,         
    breaklines=true,                 
    captionpos=b,                    
    keepspaces=true,                 
    numbers=left,                    
    numbersep=5pt,                  
    showspaces=false,                
    showstringspaces=false,
    showtabs=false,
    tabsize=4
}
\lstset{style=mystyle}


%%%%%%%%%%% MATHEMATICS %%%%%%%%%%%

% AMS packages
\usepackage{amsmath}
\usepackage{amsthm} % Better \newtheorem
\usepackage{amsfonts}
\usepackage{amssymb}
\usepackage{mathrsfs} % For \mathscr

% Scalable parentheses that behave like \mathopen and \mathclose
\usepackage{mleftright} 
\mleftright % Redefine \left as \mleft and \right as \mright

% Tag only referenced equations (this is a bad package, as it requires etextools, which is buggy and abandoned by its author)
% \expandafter\def\csname ver@etex.sty\endcsname{3000/12/31}
% \let\globcount\newcount
% \usepackage{autonum}

% For writing \overset{text}&{=} in align environment
\usepackage{aligned-overset} 

% Mathematics macros
\usepackage{cancel} % Better version of the \not command
\usepackage{polynom} % Does polynomial division for you
% \usepackage{tikz-cd} % Commutative diagrams. (I don't use this anymore, instead I use Tikz graphs: https://tikz.dev/tikz-graphs)
% \tikzcdset{every label/.append style = {font = \normalsize}}

%%% Some of my own math commands that I haven't found great packages for yet
% Bar, tilde, and hat that scales with what is under them. Basically I just want these to have consistent names
\let\mathbar\overline
\let\mathtilde\widetilde
\let\mathhat\widehat

% Main number systems
\newcommand*\naturals{\mathbb{N}}
\newcommand*\integers{\mathbb{Z}}
\newcommand*\rationals{\mathbb{Q}}
\newcommand*\reals{\mathbb{R}}
\newcommand*\complexes{\mathbb{C}}

% Some of my own shorthands for correct spacing in math environments
\def\from{\colon} % Proper spacing of colon in functions f: A → B
\newlength\mylen % Isomorphic \mapsto
\settowidth\mylen{$\longleftrightarrow$}
\newcommand{\mapsbetween}{\longleftrightarrow\kern - 0.5\mylen\vline height 1.2ex depth -0.0pt\kern0.5\mylen}
\newcommand{\suchthat}{\qq{s.th.}}
\def\definedas{\coloneqq}
\def\defines{\eqqcolon}

\newcommand*{\transpose}[1]{#1^{\mathsf{T}}}
\renewcommand*{\complement}[1]{#1^{\mathsf{C}}}
\newcommand{\conjugate}[1]{\mathbar{#1}}
% \newcommand*{\conjugate}[1]{{#1}^*}
\newcommand*{\hermitianconjugate}[1]{#1^\dag}
\newcommand*{\inverse}[1]{#1^{-1}}
\newcommand*{\dual}[1]{#1^*}

% Group action
\newcommand*\actson{\curvearrowright}


%%%%%%%%%%% MISCELLANEOUS %%%%%%%%%%%

% Cleverref automatically detects if you are referencing a figure, table, or equation etc
% This has to be loaded after babel, hyperref, amsmath, etc
\usepackage[noabbrev, nameinlink]{cleveref}
\crefname{equation}{}{}
\iflanguage{swedish}{ % Use Oxford comma
	\newcommand{\creflastconjunction}{, och\nobreakspace}
}{}
\iflanguage{english}{
	\newcommand{\creflastconjunction}{, and\nobreakspace}
}{}

% Theorem and proof environments, following the amsthm docs
% This has to be done after loading cleveref
% Beware, using package parskip will fuck with the spacing of these environments
\iflanguage{swedish}{
    \theoremstyle{plain}
    \newtheorem{theorem}{Sats}
    \newtheorem*{theorem*}{Sats}
    \newtheorem{proposition}[theorem]{Proposition}
    \newtheorem*{proposition*}{Proposition}
    \newtheorem{corollary}{Följdsats}[theorem]
    \newtheorem*{corollary*}{Följdsats}
    \newtheorem{lemma}[theorem]{Lemma}
    \newtheorem*{lemma*}{Lemma}
	\newtheorem{conjecture}[theorem]{Förmodan}
	\newtheorem*{conjecture*}{Förmodan}

    \theoremstyle{definition}
    \newtheorem{definition}[theorem]{Definition}
    \newtheorem*{definition*}{Definition}
	\newtheorem*{example}{Exempel}

    \theoremstyle{remark}
	\newtheorem*{remark}{Kommentar}
}{}
\iflanguage{english}{
    \theoremstyle{plain}
    \newtheorem{theorem}{Theorem}
    \newtheorem*{theorem*}{Theorem}
    \newtheorem{proposition}[theorem]{Proposition}
    \newtheorem*{proposition*}{Proposition}
    \newtheorem{corollary}{Corollary}[theorem]
    \newtheorem{corollary*}{Corollary}
    \newtheorem{lemma}[theorem]{Lemma}
    \newtheorem*{lemma*}{Lemma}
	\newtheorem{conjecture}[theorem]{Conjecture}
	\newtheorem{conjecture*}{Conjecture}

    \theoremstyle{definition}
    \newtheorem{definition}[theorem]{Definition}
    \newtheorem*{definition*}{Definition}
	\newtheorem*{example}{Example}

    \theoremstyle{remark}
	\newtheorem*{remark}{Remark}
}{}

% In-pdf comments through \todo command
% \setlength{\marginparwidth}{2cm} % Silence warning about margin size
% \usepackage{todonotes}
% The todonotes package is very slow, the following setup is a bit uglier but saves alot of compilation time. It may be also used with the "inline" option.
\usepackage{xstring}
\usepackage[]{marginnote}
\newcommand{\todo}[2][]{%
	\IfStrEqCase{#1}{%
        {inline}{{\color{red}#2}}%
	}[%
		\marginnote{\color{red}#2}%
	]
}

% Generate tables from csv files
\usepackage{csvsimple} 
